\chapter{Lars Hörmander --- Fields Medal 1962}\label{ch1962larshormander:lars-huxf6rmander-fields-medal-1962-a-comprehensive-report}

\begin{center}\rule{0.5\linewidth}{0.5pt}\end{center}

\textbf{Prepared:} July 2026 \textbf{Subject:} Lars Valter Hörmander (b.
January 24, 1931, Mjällby, Sweden) \textbf{Position:} Professor, Lund
University; University of Stockholm \textbf{Doctoral Advisor:} Lars
Gårding (Lund University) \textbf{Fields Medal Citation:} \emph{``Worked
in partial differential equations. Specifically, contributed to the
general theory of linear differential operators. The questions go back
to one of Hilbert's problems at the 1900 congress.''}

\begin{center}\rule{0.5\linewidth}{0.5pt}\end{center}

\section{1. Introduction}\label{ch1962larshormander:introduction}

Lars Hörmander was awarded the Fields Medal in 1962 at age 31 for his
significant contributions to the theory of linear partial
differential equations (PDEs). His work substantially reshaped the
field, unifying it through the interplay of functional analysis,
differential geometry, and Fourier analysis. He developed the modern
framework of microlocal analysis, including pseudodifferential
operators, Fourier integral operators, and the theory of wave front
sets.

\begin{prereqs} Real analysis (Lebesgue measure, $L^p$ spaces), Fourier analysis basics, basic functional analysis. \end{prereqs}

\begin{objectives} State the H\'{o}rmander condition for hypoellipticity; explain the pseudodifferential operator framework; describe the $L^2$ theory of the $\overline{\partial}$ operator. \end{objectives}


Hörmander's early work completely classified constant-coefficient
hypoelliptic operators, solving a problem that had been open since
Petrowsky in the 1930s. His geometric explanation of the Hans Lewy
nonsolvability phenomenon provided the first clear insight into why some
linear PDEs have no solutions. His \(L^2\) estimates for the
\(\bar\partial\) operator revolutionized several complex variables by
replacing intricate geometric constructions with transparent analytic
arguments. His proof that sums of squares of vector fields are
hypoelliptic under a Lie bracket spanning condition became a cornerstone
of sub-Riemannian geometry and control theory.

Hörmander's major contributions include: the complete classification of
constant-coefficient linear PDEs; \(L^2\) estimates for the
\(\bar\partial\)-Neumann problem; the theory of pseudodifferential
operators; Fourier integral operators and propagation of singularities;
the Hörmander condition for sums of squares; the sharp Weyl law for
spectral asymptotics; and the geometric theory of the Cauchy problem.
His four-volume treatise \emph{The Analysis of Linear Partial
Differential Operators} (Springer, 1983--1985) remains the definitive
reference in the field and won the 2006 Leroy P. Steele Prize for
Mathematical Exposition.

\begin{center}\rule{0.5\linewidth}{0.5pt}\end{center}

\section{2. Biographical Background}\label{ch1962larshormander:biographical-background}

\subsection{2.1 Early Life and
Education}\label{ch1962larshormander:early-life-and-education}

Lars Valter Hörmander was born on January 24, 1931, in Mjällby, a small
village in Blekinge, southern Sweden. His father was a schoolteacher who
encouraged his intellectual development. Hörmander showed exceptional
mathematical talent early on: he skipped two years of elementary school
and moved to Lund in 1946 at age 15 to attend secondary school. He
entered Lund University in 1948, where he met Lars Gårding, who would
become his doctoral advisor. Hörmander absorbed the Bourbakian style of
rigorous, systematic exposition that would characterize his own writing.

\subsection{2.2 Academic Career}\label{ch1962larshormander:academic-career}

\begin{itemize}
\tightlist
\item
 \textbf{1948--1955:} Studies at Lund University (B.Sc., Licentiate,
 Ph.D.).
\item
 \textbf{1955:} Ph.D.~under Lars Gårding. Thesis: ``On the Theory of
 General Partial Differential Operators.''
\item
 \textbf{1955--1956:} Postdoctoral visits to the University of Chicago,
 Courant Institute, and Stanford University.
\item
 \textbf{1957--1964:} Professor, University of Stockholm.
\item
 \textbf{1962:} Fields Medal at the ICM in Stockholm.
\item
 \textbf{1964--1967:} Professor, Stanford University.
\item
 \textbf{1967--1997:} Professor, Lund University.
\item
 \textbf{1988:} Wolf Prize in Mathematics.
\item
 \textbf{2006:} Leroy P. Steele Prize for Mathematical Exposition.
\item
 \textbf{2012:} Died at age 81 in Lund, Sweden.
\end{itemize}

\subsection{2.3 Doctoral Work}\label{ch1962larshormander:doctoral-work}

Hörmander's Ph.D.~thesis already contained significant results. He studied
linear PDEs with constant coefficients and gave a complete
characterization of when such operators are hypoelliptic --- that is,
when every distributional solution is smooth. A constant-coefficient
operator \(P(D) = \sum_{|\alpha| \leq m} a_\alpha D^\alpha\) is
hypoelliptic if \(P(D)u \in C^\infty\) implies \(u \in C^\infty\). This
thesis solved a problem that had been open since Petrowsky's work in the
1930s, and it immediately established Hörmander as a leading figure in
analysis. The work was published in \emph{Acta Mathematica} in 1955 and
remains a classic in the field.

\begin{center}\rule{0.5\linewidth}{0.5pt}\end{center}

\section{3. Classification of Constant-Coefficient Linear
PDEs}\label{ch1962larshormander:classification-of-constant-coefficient-linear-pdes}

\subsection{3.1 Background}\label{ch1962larshormander:background}

A linear partial differential operator with constant coefficients is an
expression \(P(D) = \sum_{|\alpha| \leq m} a_\alpha D^\alpha\), where
\(\alpha = (\alpha_1,\ldots,\alpha_n)\) is a multi-index,
\(D^\alpha = (-i\partial)^{\alpha_1} \cdots (-i\partial)^{\alpha_n}\),
and \(a_\alpha \in \mathbb{C}\). The symbol of \(P(D)\) is the
polynomial \(P(\xi) = \sum_{|\alpha| \leq m} a_\alpha \xi^\alpha\) for
\(\xi \in \mathbb{R}^n\).

A fundamental problem in PDE theory is to determine when \(P(D)u = f\)
is solvable and when its solutions are smooth. An operator is
\textbf{hypoelliptic} if \(P(D)u \in C^\infty\) implies
\(u \in C^\infty\). The Laplace operator \(\Delta\) is hypoelliptic
because elliptic regularity implies \(u \in C^\infty\) whenever
\(\Delta u \in C^\infty\). The wave operator \(\partial_t^2 - \Delta\),
by contrast, is not hypoelliptic, as weak solutions can have
singularities propagating along light cones.

\subsection{3.2 Hörmander's
Characterization}\label{ch1962larshormander:huxf6rmanders-characterization}

In his thesis, Hörmander gave a complete characterization of
hypoellipticity for constant-coefficient operators. He proved \(P(D)\)
is hypoelliptic if and only if:

\[\lim_{|\xi| \to \infty} \frac{|\partial^\alpha P(\xi)|}{|P(\xi)|^{1/|\alpha|}} = 0 \quad \text{for all } \alpha \neq 0\]

Equivalently, for every multi-index \(\alpha \neq 0\), there exist
constants \(C_\alpha, R_\alpha\) such that
\(|\partial^\alpha P(\xi)| \leq C_\alpha |P(\xi)|^{1/|\alpha|}\) for
\(|\xi| \geq R_\alpha\). This condition means the polynomial \(P\) does
not oscillate too rapidly relative to its magnitude, which guarantees
the fundamental solution (the inverse Fourier transform of \(1/P(\xi)\))
is smooth away from the origin.

\subsection{3.3 Examples and
Consequences}\label{ch1962larshormander:examples-and-consequences}

Elliptic operators satisfy \(|P(\xi)| \sim |\xi|^m\) for large
\(|\xi|\), so the condition holds and they are hypoelliptic. Parabolic
operators like the heat operator \(\partial_t - \Delta\) satisfy
\(|P(\xi)| \sim |\xi|^2\) in appropriate regions and are also
hypoelliptic. The wave operator \(\partial_t^2 - \Delta\) fails the
condition because \(\xi_1^2 - \xi_2^2 - \cdots - \xi_n^2\) can be small
relative to its derivatives along the light cone. Hörmander also proved
every constant-coefficient operator is locally solvable, using the
Malgrange--Ehrenpreis theorem. For variable-coefficient operators, this
fails dramatically --- the Hans Lewy example shows some operators are
not locally solvable anywhere.

\subsection{3.4 Impact}\label{ch1962larshormander:impact}

This work closed a chapter that began with Hilbert's 1900 Paris lecture,
where Hilbert listed the problem of developing a general theory of
linear PDEs as one of his 23 problems. Hörmander's results were the
first major step toward answering Hilbert's question. The work
established Hörmander's reputation and set the stage for his subsequent
research on variable-coefficient operators.

\begin{center}\rule{0.5\linewidth}{0.5pt}\end{center}

\section{\texorpdfstring{4. \(L^2\) Estimates for the \(\bar\partial\)
Operator}{4. L\^{}2 Estimates for the \textbackslash bar\textbackslash partial Operator}}\label{ch1962larshormander:l2-estimates-for-the-barpartial-operator}

\subsection{4.1 Background}\label{ch1962larshormander:background-1}

On a complex manifold \(M\), the \(\bar\partial\) operator maps
\((0,q)\)-forms to \((0,q+1)\)-forms and satisfies
\(\bar\partial^2 = 0\). The Dolbeault cohomology groups
\(H^{p,q}(M) = \ker\bar\partial/\text{im}\,\bar\partial\) measure the
extent to which \(\bar\partial\)-closed forms are exact. The
Grothendieck--Dolbeault lemma states that on a Stein manifold,
\(H^{p,q}(M) = 0\) for \(q \geq 1\). The core problem in several complex
variables is: given a \(\bar\partial\)-closed \((0,q)\)-form \(f\),
solve \(\bar\partial u = f\) with optimal estimates. This is essential
for constructing holomorphic functions and sections of holomorphic
bundles.

\subsection{4.2 Hörmander's Theorem}\label{ch1962larshormander:huxf6rmanders-theorem}

Hörmander proved a landmark result. Let \(\Omega \subset \mathbb{C}^n\)
be pseudoconvex and \(\phi\) a plurisubharmonic function on \(\Omega\).
For any \(\bar\partial\)-closed \((0,q)\)-form \(f\) with
\(\int_\Omega |f|^2 e^{-\phi} \, dV < \infty\), there exists \(u\) with
\(\bar\partial u = f\) and:

\[\int_\Omega |u|^2 e^{-\phi} (1+|z|^2)^{-2} dV \leq \int_\Omega |f|^2 e^{-\phi} dV\]

The estimate depends only on the plurisubharmonicity of \(\phi\), not on
the detailed geometry of \(\Omega\). The proof uses the
\textbf{Bochner--Kodaira--Nakano identity}:

\[\Box = \nabla^*\nabla + [i\Theta(E), \Lambda]\]

where \(\Theta(E)\) is the curvature of the holomorphic vector bundle
and \(\Lambda\) is adjoint to multiplication by the Kähler form. By
choosing a weight \(\phi\) with positive curvature
\(i\partial\bar\partial\phi\), Hörmander obtained a lower bound for
\(\Box\) that allowed him to solve \(\bar\partial u = f\) with optimal
regularity.

\subsection{\texorpdfstring{4.3 The \(\bar\partial\)-Neumann
Problem}{4.3 The \textbackslash bar\textbackslash partial-Neumann Problem}}\label{ch1962larshormander:the-barpartial-neumann-problem}

The \(\bar\partial\)-Neumann problem on bounded pseudoconvex domains
with smooth boundary was also solved by Hörmander's \(L^2\) estimates.
The operator
\(\Box_q = \bar\partial\bar\partial^*_q + \bar\partial^*_{q+1}\bar\partial\)
has a bounded inverse \(N_q\) on \(L^2_{(0,q)}(\Omega)\), and the
Bergman projection satisfies
\(P = I - \bar\partial^* N_1 \bar\partial\).

\subsection{4.4 Applications}\label{ch1962larshormander:applications}

Hörmander's \(L^2\) estimates have numerous applications: solving the
Levi problem (every pseudoconvex domain is a domain of holomorphy); the
excision of compact sets in Stein manifolds; construction of holomorphic
functions with prescribed zeros; and estimates for the Bergman kernel.
Bombieri and Skoda extended the method to prove effective results in
algebraic geometry.

\begin{center}\rule{0.5\linewidth}{0.5pt}\end{center}

\section{5. The Hörmander Condition for Sums of
Squares}\label{ch1962larshormander:the-huxf6rmander-condition-for-sums-of-squares}

\subsection{5.1 Statement}\label{ch1962larshormander:statement}

Consider smooth real vector fields \(X_1,\ldots,X_k\) on \(M\) and the
operator \(L = X_1^2 + \cdots + X_k^2\). This is a degenerate elliptic
operator --- its principal symbol vanishes in directions orthogonal to
the span of the vector fields. The prototype is the
\textbf{sub-Laplacian} on the Heisenberg group \(\mathbb{H}^1\):

\[L = X^2 + Y^2, \quad X = \partial_x - \frac{y}{2}\partial_t, \quad Y = \partial_y + \frac{x}{2}\partial_t\]

This operator is not elliptic (it is degenerate in the \(t\)-direction),
yet \(Lu \in C^\infty\) implies \(u \in C^\infty\) --- it is
hypoelliptic.

\subsection{5.2 Hörmander's Theorem}\label{ch1962larshormander:huxf6rmanders-theorem-1}

In his seminal 1967 paper, Hörmander proved a general theorem: if the
Lie algebra generated by \(X_1,\ldots,X_k\) spans \(T_xM\) at every
point (the \textbf{Hörmander condition}), then \(L\) is hypoelliptic.
Let \(\mathcal{L}\) be the Lie algebra generated by \(X_1,\ldots,X_k\)
(closed under commutators). If for every \(x \in M\),
\(\mathcal{L}(x) = T_xM\), then \(L\) is hypoelliptic. For the
Heisenberg group, \([X,Y] = \partial_t\), so \(X,Y,[X,Y]\) span
\(T\mathbb{R}^3\), and the condition holds.

\subsection{5.3 Geometry and Regularity}\label{ch1962larshormander:geometry-and-regularity}

The Hörmander condition defines \textbf{sub-Riemannian geometry}
(Carnot--Carathéodory geometry). A sub-Riemannian manifold consists of a
distribution \(D \subset TM\) together with a metric on \(D\), such that
Lie brackets of sections of \(D\) generate \(TM\). The associated
sub-Laplacian is the sum of squares of orthonormal vector fields framing
\(D\). The condition implies the control distance is well-defined and
the sub-Laplacian satisfies scale-invariant parabolic Harnack
inequalities.

Hörmander's proof combined Fourier analysis on nilpotent Lie groups, the
Rothschild--Stein lifting theorem, and the construction of
pseudodifferential parametrices. The heat kernel for \(L\) has
small-time asymptotics involving the sub-Riemannian geodesic distance.

\subsection{5.4 Impact}\label{ch1962larshormander:impact-1}

The Hörmander condition is now essential in control theory (the Lie
bracket condition for controllability), hypoelliptic operators on Lie
groups (the Rockland conjecture, proved by Helffer and Nourrigat),
sub-Riemannian geometry, and hypoelliptic diffusion.

\begin{center}\rule{0.5\linewidth}{0.5pt}\end{center}

\section{6. Pseudodifferential
Operators}\label{ch1962larshormander:pseudodifferential-operators}

\subsection{6.1 Symbols and Classes}\label{ch1962larshormander:symbols-and-classes}

A pseudodifferential operator (\(\Psi\)DO) is:

\[(Au)(x) = \frac{1}{(2\pi)^n} \int e^{i x\cdot\xi} a(x,\xi) \hat u(\xi) \, d\xi\]

For a differential operator
\(P(x,D) = \sum_{|\alpha| \leq m} a_\alpha(x) D^\alpha\), the symbol is
\(p(x,\xi) = \sum_{|\alpha| \leq m} a_\alpha(x) \xi^\alpha\).

Hörmander introduced symbol classes \(S^m_{\rho,\delta}\):

\[|\partial_\xi^\alpha \partial_x^\beta a(x,\xi)| \leq C_{\alpha\beta} (1+|\xi|)^{m-\rho|\alpha|+\delta|\beta|}\]

with \(0 \leq \delta < \rho \leq 1\). The standard case
\(\rho=1,\delta=0\) gives classical symbols.

\subsection{6.2 The Calculus}\label{ch1962larshormander:the-calculus}

\begin{itemize}
\item
 \textbf{Composition}: \(A \circ B \in \Psi^{m_1+m_2}\) with symbol:

 \[a \circ b \sim \sum_{\alpha} \frac{(-i)^{|\alpha|}}{\alpha!} \partial_\xi^\alpha a \cdot \partial_x^\alpha b\]
\item
 \textbf{Adjoints}: \(A^* \in \Psi^m\) with symbol given by an
 asymptotic series.
\item
 \textbf{Ellipticity}: If the symbol is invertible for large \(|\xi|\),
 the operator has a parametrix.
\item
 \textbf{Sharp Gårding inequality}:
 \((Au,u) \geq -C\|u\|^2_{H^{(m-1)/2}}\) for \(a \geq 0\).
\end{itemize}

\subsection{6.3 Impact}\label{ch1962larshormander:impact-2}

The \(\Psi\)DO calculus is essential for spectral geometry, the
Atiyah--Singer index theorem, and geometric analysis. Hörmander's Volume
III is the definitive reference.

\begin{center}\rule{0.5\linewidth}{0.5pt}\end{center}

\section{7. Fourier Integral
Operators}\label{ch1962larshormander:fourier-integral-operators}

\subsection{7.1 Definition}\label{ch1962larshormander:definition}

An FIO generalizes \(\Psi\)DOs for hyperbolic equations:

\[(Ku)(x) = \frac{1}{(2\pi)^n} \int e^{i\phi(x,\xi)} a(x,\xi) \hat u(\xi) \, d\xi\]

where \(\phi\) is a non-degenerate phase homogeneous of degree 1 in
\(\xi\).

\subsection{7.2 Hörmander's Theory}\label{ch1962larshormander:huxf6rmanders-theory}

Hörmander (1968--1971) established: - FIOs correspond to canonical
relations \(\Lambda \subset T^*X \times T^*Y\) (Lagrangian
submanifolds). - The principal symbol transforms as a half-density on
\(\Lambda\). - Composition corresponds to composition of canonical
relations.

\subsection{7.3 Propagation of
Singularities}\label{ch1962larshormander:propagation-of-singularities}

With Duistermaat, Hörmander proved that for \(P\) of real principal
type, \(\text{WF}(u)\) propagates along bicharacteristics:

\[\dot x = \partial_\xi p, \quad \dot \xi = -\partial_x p, \quad p(x,\xi)=0\]

This is the foundation of microlocal analysis, with applications in
seismic imaging and tomography.

\begin{center}\rule{0.5\linewidth}{0.5pt}\end{center}

\section{8. The Cauchy Problem and
Pseudoconvexity}\label{ch1962larshormander:the-cauchy-problem-and-pseudoconvexity}

\subsection{8.1 The Hans Lewy Example}\label{ch1962larshormander:the-hans-lewy-example}

In 1957, Hans Lewy discovered a shocking phenomenon: there exists a
smooth linear PDE with no local solutions anywhere. The operator
\(L u = \frac{\partial u}{\partial \bar z} - i\bar z \frac{\partial u}{\partial t}\)
has the property that for most \(f \in C^\infty\), \(L u = f\) has no
distributional solution in any open set.

Hörmander explained this geometrically. For \(L = X + iY\), local
solvability requires \([X,Y]\) to be in the span of \(X\) and \(Y\) (the
Levi condition). For the Lewy operator, \([X,Y] = \partial_t\) is
transverse.

\subsection{8.2 Pseudoconvexity and Carleman
Estimates}\label{ch1962larshormander:pseudoconvexity-and-carleman-estimates}

Hörmander proved that Carleman estimates
\(\|u e^{\tau\phi}\| \leq C\|Pu e^{\tau\phi}\|\) exist iff:

\[\{p,\phi\}=0,\; p=0 \implies \{\overline{p},\{p,\phi\}\} > 0\]

This ensures uniqueness in the Cauchy problem.

\begin{center}\rule{0.5\linewidth}{0.5pt}\end{center}

\section{9. Spectral Asymptotics and the Weyl
Law}\label{ch1962larshormander:spectral-asymptotics-and-the-weyl-law}

\subsection{9.1 The Sharp Weyl Law}\label{ch1962larshormander:the-sharp-weyl-law}

Using FIOs, Hörmander proved:

\[N(\lambda) = \frac{\omega_n}{(2\pi)^n} \text{Vol}(M) \lambda^{n/2} + O(\lambda^{(n-1)/2})\]

The proof uses the wave kernel \(\cos(t\sqrt{-\Delta})\) (an FIO for
small \(t\)) and stationary phase. The remainder is optimal, as shown by
\(S^n\).

\subsection{9.2 The Poisson Relation}\label{ch1962larshormander:the-poisson-relation}

The singularities of \(\text{Tr}\cos(t\sqrt{-\Delta})\) are contained in
the length spectrum of closed geodesics, fundamental in quantum chaos
and inverse spectral theory.

\begin{center}\rule{0.5\linewidth}{0.5pt}\end{center}

\section{10. The Weyl Calculus and Hörmander
Metrics}\label{ch1962larshormander:the-weyl-calculus-and-huxf6rmander-metrics}

\subsection{10.1 Weyl Quantization}\label{ch1962larshormander:weyl-quantization}

\[(A^w u)(x) = \frac{1}{(2\pi)^n} \iint e^{i\langle x-y,\xi\rangle}
a\left(\frac{x+y}{2},\xi\right) u(y) \, dy\, d\xi\]

transforms naturally under symplectic maps, ideal for operators with
double characteristics.

\subsection{10.2 Hörmander Metrics}\label{ch1962larshormander:huxf6rmander-metrics}

Metrics \(g_{x,\xi}\) on \(T^*X\) satisfying the uncertainty principle,
slow variation, and temperateness define symbol classes \(S(m,g)\). They
support the sharp Gårding and Hörmander--Melin inequalities.

\begin{center}\rule{0.5\linewidth}{0.5pt}\end{center}

\section{11. The Four-Volume Treatise}\label{ch1962larshormander:the-four-volume-treatise}

\emph{The Analysis of Linear Partial Differential Operators} (Springer,
1983--1985) won the 2006 Steele Prize: - \textbf{Vol. I}: Distribution
theory, Fourier analysis, Sobolev spaces. - \textbf{Vol. II}:
Constant-coefficient operators --- classification, hypoellipticity,
fundamental solutions. - \textbf{Vol. III}: Pseudodifferential operators
--- calculus, elliptic regularity, spectral theory. - \textbf{Vol. IV}:
Fourier integral operators --- propagation, Cauchy problem, spectral
asymptotics.

It is the definitive reference, self-contained and rigorous.

\begin{center}\rule{0.5\linewidth}{0.5pt}\end{center}

\section{12. Hyperfunctions and the Edge of the
Wedge}\label{ch1962larshormander:hyperfunctions-and-the-edge-of-the-wedge}

\textbf{Hyperfunctions} (Sato) generalize distributions via boundary
values of holomorphic functions. Hörmander gave a Fourier-analytic
treatment: every compactly supported hyperfunction is the Fourier
transform of an entire function of exponential type.

The \textbf{edge of the wedge theorem}: let
\(\Gamma_1,\Gamma_2 \subset \mathbb{R}^n\) be open convex cones, \(f_j\)
holomorphic on \(\Omega+i\Gamma_j\) with equal boundary values.
Hörmander proved they extend to a holomorphic function on the convex
hull, using Fourier analysis.

\begin{center}\rule{0.5\linewidth}{0.5pt}\end{center}

\section{13. The Radon Transform and Integral
Geometry}\label{ch1962larshormander:the-radon-transform-and-integral-geometry}

The Radon transform
\(Rf(\omega,p) = \int_{x\cdot\omega=p} f(x)\, dS(x)\) is an FIO with
canonical relation:

\[\Lambda = \{((\omega,p), (x,\xi)): x\cdot\omega = p, \xi = |\xi|\omega\}\]

\(R^*R\) is an elliptic \(\Psi\)DO. For the X-ray transform,
\(R^*R = c_n(-\Delta)^{-1/2}\). This provides the foundation for CT
scanning.

\begin{center}\rule{0.5\linewidth}{0.5pt}\end{center}

\section{\texorpdfstring{14. CR Geometry and the \(\bar\partial_b\)
Complex}{14. CR Geometry and the \textbackslash bar\textbackslash partial\_b Complex}}\label{ch1962larshormander:cr-geometry-and-the-barpartial_b-complex}

A CR structure on \(M^{2n+1}\) is
\(T^{1,0}M \subset TM\otimes\mathbb{C}\) of rank \(n\) with
\(T^{1,0}M \cap \overline{T^{1,0}M} = \{0\}\) and
\([T^{1,0}M, T^{1,0}M] \subset T^{1,0}M\). Hörmander studied
\(\Box_b = \bar\partial_b\bar\partial_b^* + \bar\partial_b^*\bar\partial_b\),
proving subelliptic estimates when the Levi form has sufficient
positivity.

\begin{center}\rule{0.5\linewidth}{0.5pt}\end{center}

\section{15. Operators with Double
Characteristics}\label{ch1962larshormander:operators-with-double-characteristics}

\(P\) has double characteristics if \(p_m\) vanishes to order 2 on a
symplectic submanifold \(\Sigma\). Examples: magnetic Schrödinger and
Klein--Gordon operators. Hörmander reduced such operators to a model on
a symplectic vector space. The subprincipal symbol determines
hypoellipticity: if the Hamilton map satisfies the
Ivrii--Petkov--Hörmander condition, \(P\) is hypoelliptic with loss of
\(1/2\) derivative.

\begin{center}\rule{0.5\linewidth}{0.5pt}\end{center}

\section{16. Gevrey Classes}\label{ch1962larshormander:gevrey-classes}

\(f \in G^s(\Omega)\) if
\(|\partial^\alpha f(x)| \leq C^{|\alpha|+1}(|\alpha|!)^s\) on compact
sets, interpolating \(C^\infty\) (\(s=1\)) and real-analytic
(\(s=\infty\)). Hörmander proved the Cauchy problem for weakly
hyperbolic equations is well-posed in \(G^s\) for \(s\) sufficiently
small.

\begin{center}\rule{0.5\linewidth}{0.5pt}\end{center}

\section{17. Hyperbolic Equations and Energy
Estimates}\label{ch1962larshormander:hyperbolic-equations-and-energy-estimates}

\(P\) is strictly hyperbolic wrt \(t\) if \(p_m(x,\xi)\) has \(m\)
distinct real roots. Hörmander proved:

\[\sum_{j=0}^{m-1} \|u\|_{C^j([0,T]; H^{s-j})} \leq C_T \left(
\sum_{j=0}^{m-1} \|u_j\|_{H^{s-j}} + \int_0^T \|f\|_{H^s} dt\right)\]

The \textbf{Mizohata condition}: the Cauchy problem is \(C^\infty\)
well-posed iff \(p_m(x,\xi+i\tau d\phi) \neq 0\) for
\(\tau \geq \tau_0\).

\begin{center}\rule{0.5\linewidth}{0.5pt}\end{center}

\section{18. Theory of Parametrices}\label{ch1962larshormander:theory-of-parametrices}

A parametrix for \(P\) is \(Q\) with \(PQ = I+R\), \(QP = I+S\) (\(R,S\)
smoothing). Hörmander constructed parametrices: - For elliptic
operators: as \(\Psi\)DOs. - For operators of real principal type: as
FIOs. - For the wave equation: \(\cos(t\sqrt{-\Delta}) = \frac{1}{2\pi}
\int e^{i(x-y)\xi} \cos(t|\xi|) d\xi + \text{lower terms}\).

\begin{center}\rule{0.5\linewidth}{0.5pt}\end{center}

\section{19. The Wave Front Set}\label{ch1962larshormander:the-wave-front-set}

\((x_0,\xi_0) \notin \text{WF}(u)\) if there exists \(A \in \Psi^0\)
with symbol 1 near \((x_0,\xi_0)\) such that \(Au \in C^\infty\).
Hörmander proved: - \(\text{WF}(Au) = \text{WF}(u)\) for elliptic \(A\).
- For \(Pu = f\) of real principal type, \(\text{WF}(u)\) propagates
along bicharacteristics. - The analytic wave front set uses the FBI
transform.

\begin{center}\rule{0.5\linewidth}{0.5pt}\end{center}

\section{20. The Hörmander--Melin
Inequality}\label{ch1962larshormander:the-huxf6rmandermelin-inequality}

When \(a \geq 0\) vanishes on a symplectic submanifold:

\[(A^w u, u) \geq -C_1 \|u\|^2_{H^{(m-1)/2}} -
C_2 \|u\|^2_{H^{(m-1)/2}} \log\left(2 + \frac{\|u\|^2_{H^{m/2}}}
{\|u\|^2_{H^{(m-1)/2}}}\right)\]

Essential for subelliptic operators and \(\Box_b\) on CR manifolds.

\begin{center}\rule{0.5\linewidth}{0.5pt}\end{center}

\section{21. Awards and Recognition}\label{ch1962larshormander:awards-and-recognition}

\textbf{1962:} - Fields Medal (IMU) --- at the ICM in Stockholm.

\textbf{1988:} - Wolf Prize in Mathematics.

\textbf{2006:} - Leroy P. Steele Prize for Mathematical Exposition
(AMS).

\textbf{Elected to:} - Royal Swedish Academy of Sciences (1968). -
National Academy of Sciences (USA) (1976). - Royal Society, Foreign
Member (1980). - French Academy of Sciences, Foreign Associate (1984). -
Finnish Society of Sciences, Royal Danish Academy, Norwegian Academy.

\begin{center}\rule{0.5\linewidth}{0.5pt}\end{center}

\section{22. Selected Major
Publications}\label{ch1962larshormander:selected-major-publications}

\begin{enumerate}
\def\labelenumi{\arabic{enumi}.}
\item
 \textbf{``On the theory of general partial differential operators.''}
 \emph{Acta Mathematica} 94, 161--248, 1955. --- \emph{Thesis:
 classification of constant-coefficient hypoelliptic operators.}
\item
 \textbf{``\(L^2\) estimates and existence theorems for the
 \(\bar\partial\) operator.''} \emph{Acta Mathematica} 113, 89--152,
 1965. --- \emph{Weighted \(L^2\) estimates for \(\bar\partial\).}
\item
 \textbf{``Hypoelliptic second order differential equations.''}
 \emph{Acta Mathematica} 119, 147--171, 1967. --- \emph{The Hörmander
 condition for sums of squares.}
\item
 \textbf{``Fourier integral operators I.''} \emph{Acta Mathematica}
 127, 79--183, 1971. --- \emph{Foundational paper on FIOs.}
\item
 \textbf{``On the singularities of solutions of partial differential
 equations.''} (with J. J. Duistermaat) \emph{Comm. Pure Appl. Math.}
 25, 207--236, 1972. --- \emph{Propagation of singularities.}
\item
 \textbf{``The spectral function of an elliptic operator.''} \emph{Acta
 Mathematica} 121, 173--218, 1968. --- \emph{Sharp Weyl law.}
\item
 \textbf{``Uniqueness theorems and wave front sets.''} \emph{Arkiv för
 Matematik} 6, 1--31, 1965. --- \emph{Wave front set theory.}
\item
 \textbf{``The analysis of linear partial differential operators
 I--IV.''} \emph{Grundlehren} 256, 257, 274, 275, Springer, 1983--1985.
 --- \emph{The definitive treatise.}
\item
 \textbf{``Linear partial differential operators.''} \emph{Grundlehren}
 116, Springer, 1963. --- \emph{First comprehensive book on linear
 PDEs.}
\item
 \textbf{``An introduction to complex analysis in several variables.''}
 \emph{North-Holland}, 1966. --- \emph{Classic text using \(L^2\)
 methods.}
\item
 \textbf{``Notions of convexity.''} \emph{Birkhäuser}, 1994. ---
 \emph{Convexity in PDEs.}
\item
 \textbf{``The Weyl calculus of pseudodifferential operators.''}
 \emph{Comm. Pure Appl. Math.} 32, 359--443, 1979. --- \emph{Weyl
 calculus and Hörmander metrics.}
\end{enumerate}

\begin{center}\rule{0.5\linewidth}{0.5pt}\end{center}

\begin{center}\rule{0.5\linewidth}{0.5pt}\end{center}

\section*{Exercises}
\addcontentsline{toc}{section}{Exercises}

\begin{exercise} Verify that the Laplacian $\\Delta$ satisfies the H\\'{o}rmander condition $\\sum |a_{ij}(x)| \\leq C$. \end{exercise}

\begin{exercise} Compute the Fourier transform of $e^{-|\\xi|^2}$ and verify the Parseval identity. \end{exercise}

\begin{exercise} Show that $P(D) = -\\Delta + 1$ is elliptic by checking the principal symbol. \end{exercise}

\begin{exercise} Prove that a distribution $u$ satisfying $Pu = f$ with $P$ hypoelliptic is smooth wherever $f$ is smooth. \end{exercise}

\begin{center}\rule{0.5\linewidth}{0.5pt}\end{center}

\section*{Summary}
\addcontentsline{toc}{section}{Summary}

H\'{o}rmander developed the modern $L^2$ theory of PDEs, introducing pseudodifferential operators and the Fourier integral operator framework. His characterization of hypoellipticity and his work on the $\overline{\partial}$-Neumann problem established the functional-analytic foundations of modern PDE theory.

\section{23. References and Further
Reading}\label{ch1962larshormander:references-and-further-reading}

\subsection{Official Sources}\label{ch1962larshormander:official-sources}

\begin{itemize}
\tightlist
\item
 IMU Fields Medal 1962:
 https://www.mathunion.org/imu-awards/fields-medal/fields-medals-1962
\item
 Wolf Prize: https://wolffund.org.il/lars-hormander/
\end{itemize}

\subsection{Key Technical References}\label{ch1962larshormander:key-technical-references}

\begin{itemize}
\tightlist
\item
 Hörmander, L. \emph{The Analysis of Linear Partial Differential
 Operators I--IV}. Springer, 1983--1985.
\item
 Hörmander, L. \emph{Linear Partial Differential Operators}. Springer,
 1963.
\item
 Hörmander, L. \emph{An Introduction to Complex Analysis in Several
 Variables}. North-Holland, 1966.
\item
 Hörmander, L. \emph{Notions of Convexity}. Birkhäuser, 1994.
\end{itemize}

\subsection{Textbooks and Surveys}\label{ch1962larshormander:textbooks-and-surveys}

\begin{itemize}
\tightlist
\item
 Treves, F. \emph{Introduction to Pseudodifferential and Fourier
 Integral Operators I--II}. Plenum, 1980.
\item
 Taylor, M. E. \emph{Pseudodifferential Operators}. Princeton
 University Press, 1981.
\item
 Duistermaat, J. J. \emph{Fourier Integral Operators}. Birkhäuser,
 1996.
\item
 Grigis, A. and Sjöstrand, J. \emph{Microlocal Analysis for
 Differential Operators}. Cambridge, 1994.
\end{itemize}

\subsection{Biography and Exposition}\label{ch1962larshormander:biography-and-exposition}

\begin{itemize}
\tightlist
\item
 Gårding, L. ``The Work of Lars Hörmander.'' ICM 1962 Laudatio.
\item
 Boman, J. et al.~``Lars Hörmander (1931--2012).'' \emph{Notices AMS}
 62(4), 2015.
\item
 Atiyah, M. ``Lars Hörmander.'' \emph{Biographical Memoirs of the Royal
 Society} 59, 2013.
\end{itemize}

\begin{center}\rule{0.5\linewidth}{0.5pt}\end{center}


\begin{center}\rule{0.5\linewidth}{0.5pt}\end{center}

\section*{See Also}
\addcontentsline{toc}{section}{See Also}

This chapter relates to the following work by other Fields Medalists:
\begin{itemize}
  \item \textbf{\hyperref[ch1950laurentschwartz:laurent-schwartz-fields-medal-1950-a-comprehensive-report]{See Chapter on related work}} --- Hormander extended Schwartz's distribution theory.
  \item \textbf{\hyperref[ch2006terencetao:terence-tao-fields-medal-2006-a-comprehensive-report-on-his-work-and-impact]{See Chapter on related work}} --- Tao used Hormander's work on oscillatory integrals.
\end{itemize}

\emph{(End of report)}
